% year: תש"פ
% type: מבחן
% semester: א
% moed: א
% part: ב
% number: 3
% points: 9
% categories: derivatives, multiple-choice
% source: exams/מבחנים/תשפ/מועד א תשפ.pdf

\begin{question}
הנגזרת $y'(0)$ של הפונקציה הסתומה $y(x)$ המוגדרת על ידי המשוואה $\cos\left(\frac{\pi}{2}+x+y\right)-y=0$ בסביבה של הנקודה $(0,0)$ שווה ל-
\begin{enumerate}
  \item[(א)] $\frac12$
  \item[(ב)] $0$
  \item[(ג)] $\frac{\pi}{2}$
  \item[(ד)] $1+\frac{\pi}{2}$
  \item[(ה)] $-\frac12$
\end{enumerate}
\end{question}

\begin{hint}
גזרו את שני אגפי המשוואה לפי $x$, כאשר $y=y(x)$ (כלל השרשרת), והציבו $x=0$, $y=0$.
\end{hint}

\begin{solution}
\textbf{התשובה הנכונה: (ה).}

ראשית, הנקודה $(0,0)$ מקיימת את המשוואה: $\cos\frac{\pi}{2}-0=0$. נגזור את שני אגפי המשוואה לפי $x$, כאשר $y=y(x)$, לפי כלל השרשרת:
\[ -\sin\left(\frac{\pi}{2}+x+y\right)\cdot\left(1+y'\right) - y' = 0. \]
נציב $x=0$, $y(0)=0$; אז $\sin\frac{\pi}{2}=1$:
\[ -(1+y'(0)) - y'(0) = 0 \quad\Longrightarrow\quad -1-2y'(0)=0 \quad\Longrightarrow\quad y'(0) = -\frac12. \]
(לחלופין: $\cos\left(\frac\pi2+t\right)=-\sin t$, ולכן המשוואה היא $\sin(x+y)+y=0$; גזירה נותנת $\cos(x+y)(1+y')+y'=0$, וב-$(0,0)$: $1+2y'=0$.)
\end{solution}
